% META: {"id":"1SPE-DERIVATION-LOCAL-FR-R1","chapitre":"1SPE-DERIVATION-LOCAL","type_objet":"remediation","capacites_codes":["R1"],"status":"generated"}

%---------------------------------------------------------------
% FICHE DE REMISE À NIVEAU — R1 : Calculer une image, lire une image sur une courbe
%---------------------------------------------------------------

\subsection*{Rappel — Images et courbe représentative}

\textbf{Calculer une image} : si $f(x) = 2x^2 - 3x + 1$, alors l'image de $4$ par $f$ est
\[
  f(4) = 2 \times 4^2 - 3 \times 4 + 1 = 32 - 12 + 1 = 21.
\]

\textbf{Lire une image sur une courbe} : sur la courbe $\mathcal{C}_f$, le point d'abscisse $a$ a pour ordonnée $f(a)$. On lit $f(a)$ en projetant horizontalement sur l'axe des ordonnées.

%---------------------------------------------------------------

\begin{exercice}{1SPE-DERLOCAL-FR-R1-EX1}{1}{5}
Soit $f(x) = x^2 + 2x - 3$. Calculer les images suivantes.
\begin{enumerate}
  \item $f(0)$
  \item $f(2)$
  \item $f(-1)$
  \item $f(a)$ pour $a = 5$.
\end{enumerate}
\end{exercice}

% BEGIN-VERIFY
% from sympy import symbols
% x = symbols('x')
% f = x**2 + 2*x - 3
% assert f.subs(x, 0) == -3
% assert f.subs(x, 2) == 5
% assert f.subs(x, -1) == -4
% assert f.subs(x, 5) == 32
% END-VERIFY

\begin{corrige}{1SPE-DERLOCAL-FR-R1-EX1}

\commentaireMarge{On remplace $x$ par la valeur indiquée et on effectue les calculs.}

\textbf{Question 1.}
\[
f(0) = 0^2 + 2 \times 0 - 3 = -3.
\]

\textbf{Question 2.}
\[
f(2) = 2^2 + 2 \times 2 - 3 = 4 + 4 - 3 = 5.
\]

\textbf{Question 3.}
\[
f(-1) = (-1)^2 + 2 \times (-1) - 3 = 1 - 2 - 3 = -4.
\]

\textbf{Question 4.}
\[
f(5) = 5^2 + 2 \times 5 - 3 = 25 + 10 - 3 = 32.
\]

\end{corrige}

%---------------------------------------------------------------

\begin{exercice}{1SPE-DERLOCAL-FR-R1-EX2}{1}{5}
Soit $g(x) = \dfrac{1}{x+1}$ définie pour $x \neq -1$.
\begin{enumerate}
  \item Calculer $g(0)$, $g(1)$ et $g(3)$.
  \item Calculer $g(a)$ puis $g(a+h)$ pour des valeurs génériques $a$ et $h$.
  \item Pourquoi $g(-1)$ n'existe-t-elle pas ?
\end{enumerate}
\end{exercice}

% BEGIN-VERIFY
% from sympy import symbols, Rational
% x, a, h = symbols('x a h')
% g = 1/(x+1)
% assert g.subs(x, 0) == 1
% assert g.subs(x, 1) == Rational(1, 2)
% assert g.subs(x, 3) == Rational(1, 4)
% END-VERIFY

\begin{corrige}{1SPE-DERLOCAL-FR-R1-EX2}

\textbf{Question 1.}
\[
g(0) = \frac{1}{0+1} = 1, \qquad
g(1) = \frac{1}{1+1} = \frac{1}{2}, \qquad
g(3) = \frac{1}{3+1} = \frac{1}{4}.
\]

\commentaireMarge{Pour $g(a+h)$, on remplace $x$ par $a+h$ dans la formule.}

\textbf{Question 2.}
\[
g(a) = \frac{1}{a+1}, \qquad
g(a+h) = \frac{1}{(a+h)+1} = \frac{1}{a+h+1}.
\]

\textbf{Question 3.} Si $x = -1$, le dénominateur vaut $-1+1 = 0$. La division par zéro est interdite, donc $g(-1)$ n'existe pas.

\end{corrige}

%---------------------------------------------------------------

\begin{exercice}{1SPE-DERLOCAL-FR-R1-EX3}{1}{6}
Soit $f(x) = -x^2 + 4x$.
\begin{enumerate}
  \item Calculer $f(1)$ et $f(3)$.
  \item Calculer $f(1+h)$ en fonction de $h$ (développer et réduire).
  \item En déduire $f(1+h) - f(1)$.
\end{enumerate}
\end{exercice}

% BEGIN-VERIFY
% from sympy import symbols, expand
% x, h = symbols('x h')
% f = -x**2 + 4*x
% assert f.subs(x, 1) == 3
% assert f.subs(x, 3) == 3
% f1h = expand(f.subs(x, 1+h))
% assert f1h == -h**2 + 2*h + 3
% assert expand(f1h - 3) == -h**2 + 2*h
% END-VERIFY

\begin{corrige}{1SPE-DERLOCAL-FR-R1-EX3}

\textbf{Question 1.}
\[
f(1) = -1^2 + 4 \times 1 = -1 + 4 = 3, \qquad
f(3) = -3^2 + 4 \times 3 = -9 + 12 = 3.
\]

\commentaireMarge{On remplace $x$ par $1+h$ puis on développe $(1+h)^2 = 1 + 2h + h^2$.}

\textbf{Question 2.}
\begin{align*}
f(1+h) &= -(1+h)^2 + 4(1+h) \\
        &= -(1 + 2h + h^2) + 4 + 4h \\
        &= -1 - 2h - h^2 + 4 + 4h \\
        &= -h^2 + 2h + 3.
\end{align*}

\textbf{Question 3.}
\[
f(1+h) - f(1) = (-h^2 + 2h + 3) - 3 = -h^2 + 2h = h(-h + 2).
\]

\end{corrige}
